\documentclass{article}
\usepackage{amsmath,amssymb,hyperref}
\title{Which opportunities a sealed recovery interface has to distinguish}
\date{26 November 2026}
\begin{document}
\maketitle
\paragraph{Result.}
Weak fairness and strong fairness can demand radically different retained
information. For one named outstanding duty and \(k\ge2\) strong-fairness
pairs there is a sealed two-port family in which all modules give the same
weak-fairness summary, but an exact test against all admitted neighbors
needs at least \(\binom{k}{\lfloor k/2\rfloor}\) bits. This is tight for the
specified family. An exact reusable procedure for general finite installed
graphs works by compiling a separate suffix chart per set of disappearing
opportunities. The bound constrains a self-contained reusable interface,
not time to solve Streett emptiness on a full explicit graph.

\paragraph{Pinned semantics.}
Use the \href{https://github.com/eternal-november/eos-research-programmes/pull/617}{prior
fixed-policy monitor}: the graph stops at actual terminal payment of a
named duty. Each real unpaid transition is labeled by two finite masks
\(E,T\subseteq[k]\), \(T\subseteq E\). Its opportunities are established
in the original installed full-state context, before projection, and \(T\)
marks actual occurrences. General transition predicates are allowed; for
ordinary state enablement, labels on outgoing transitions use the same
source-state \(E\). Some opportunities may have their occurrence available
only on a payment exit, so they still label the unpaid step out of that
state. This is not a model of how to infer eligibility from a status bit
or how to obtain physical settlement. Multiple occurrences may be one
primitive joint event; the test family can also serialize them at the
expense of extra finite internal states.

An infinite unpaid path is admitted when, for each \(i\), infinite
recurrence of \(i\in E\) implies infinite recurrence of \(i\in T\).
A maximal finite unpaid path stops only with no real outgoing transition,
including no available paid termination. Count such paths, fair unpaid
paths, possible payment, and stranded prefixes separately. A stranded
unpaid prefix has no admitted maximal extension, paid or unpaid.
A successful admissibility check never silently deletes a stranded prefix.
Ports carry the whole state relevant to continuations; modules are
installed under one coherent source law. Every cross edge enters and
leaves at known ports; all initial cuts are ports. An update of a bridge
can alter enablement on other transitions incident to the same state.
The old labels or summaries are reusable only after revalidating these
source dependencies. No cross-version independence or implicit
authorship of fairness is assumed.

\paragraph{Exact finite suffix adaptation.}
For each \(J\subseteq[k]\), forbid every suffix edge with
\(E\cap J\ne\varnothing\) and require repeated visits to each occurrence
color \(i\notin J\). On an admitted strong-fair infinite path take \(J\) to
be all opportunities occurring only finitely often. There is a common
finite cut after all their last appearances, and the path from then on has
each required color infinitely often. Conversely an original prefix and
any suffix passing one such chart satisfy the strong rule. Prefix
reachability is therefore always in the unfiltered graph. No post-filter
deadend acquires the status of a real stopping point. Recomputing
opportunity from the retained edges would falsely clear opportunities
whose only actual execution leads to payment or is available before the
selected suffix cut.

Use the existing nonempty port-walk and positive color summary
\emph{after} filtering in each local body. Assemble the resulting charts
over bridges filtered by the same \(J\). Reachable cyclic strongly
connected components in each chart with all required colors give the
cross-module recurrence seeds. Repeat tours through a common known port
to realize different colors on different concrete traversals; the
endpoint must preserve the state enabling another such traversal.
A suffix eventually staying inside one local body is handled separately:
mark original ingress ports from which the unfiltered local prefix can
reach a filtered internal accepting component for any \(J\).
These two cases cover every infinite run, according to whether it
crosses modules infinitely or only finitely often.
Each accepting component has a nonempty cycle, even if its required
color set is empty. The quotient verdict does not produce a continuous
physical motion.

Finite stop/paid status and extension use the original chart.
From the fair seeds, paid exits and real stops, propagate productivity
backwards on original port reachability. For closed interior traps use
original local bottom components that contain neither a port, a real stop,
a paid exit, nor a vertex in any filtered fair component. Mark original
ports from which such a bottom is reachable. Exactly as in the weak
case, an unproductive reachable port or a reachable marked bottom makes
a stranded prefix. For sufficiency, every local interior state can reach
an original bottom component; either it reaches a port with a possible
extension, reaches a paid exit or finite stop, or reaches a filtered
fair component. Conversely a bad closed interior bottom has none of those
options. This proves extension as well as detecting a fair unpaid loop.

At \(p\) declared ports, full charts store
\(O(2^k p^2(k+1)+p)\) bits beyond fixed names/version metadata for this
verdict. Compatibility and potentially large construction and validation
costs are not counted as zero. To decide on the full graph, compare
against standard \href{https://hdl.handle.net/1813/7218}{Streett nonemptiness}
and the \href{https://lmcs.episciences.org/3953/pdf}{good-component
algorithms of Chatterjee, Henzinger and Loitzenbauer}. Our simple direct
comparison prunes opportunity edges whose matching occurrence is absent
inside a candidate component, then repeats the SCC check on the remainder.
No new bound for their graph problem is claimed. Subdividing edges into
labeled vertices gives the usual vertex formulation. The use here is
reuse of frozen sealed portions, not superiority to equal factoring.

\paragraph{Sharp contextual information bound.}
Set \(d=\lfloor k/2\rfloor\) and let \(\mathcal A\) be the \(d\)-element
subsets of \([k]\). A module is determined by an arbitrary
\(S\subseteq\mathcal A\). From entry \(p\) to exit \(q\), it allows
one excursion per \(A\in S\), plus a fixed fallback excursion labeled
by the full opportunity mask. On each excursion the first edge is
unlabeled and the second makes exactly its specified \(A\) (or the full
set) available but executes none of them. At that second edge's source,
those choices are genuinely offered as alternative terminal payment
transitions; none lies inside the unpaid monitor. Supply a separate
unconditional payment option at every cut, so there is no finite
unpaid stop or stranded prefix. These are abstract qualified transitions,
not asserted functions of the simulator in the native eligibility study.

Neighbor \(K\in\mathcal A\) returns \(q\) to \(p\) through a second
module, executing precisely \(K\) on each return. The occurrences are
enabled in the return's actual context. There are no other unpaid
cycles or occurrences. An unpaid infinite run must repeat excursions.
If \(K\in S\), repeat its own excursion and the return; the only recurring
opportunities are executed on each cycle. If \(K\notin S\), every
excursion hits \([k]\setminus K\) because no distinct equal-cardinality
set lies inside \(K\), and the full fallback also hits it.
By the finite pigeonhole principle some unexecuted opportunity recurs
infinitely, so no strongly fair unpaid run exists. This is a
constructive decoder: the neighbor's fair-unpaid verdict reads exactly
the bit \(K\in S\). Payment feasibility and prefix extension stay
constant across these comparisons.

The ordinary local weak-fairness summary is independent of \(S\):
the unlabeled first step supplies every weak color
\(\neg E_i\lor T_i\), there is just one nonempty boundary pair, payment
and actual-outgoing flags agree, and none of the bodies has a local
cycle. The full fallback also makes the union of \(E\) across the
pair equal for every \(S\); the union of \(T\) is empty. An interface
cannot recover the lost distinctions by adding separate marginals.

Now any deterministic fixed-width summary whose decoder sees only
it and the chosen neighbor (not the hidden body through a registry or
query) must distinguish all \(2^{|\mathcal A|}\) choices. If two summaries
coincide, select a \(K\) on which the two choices differ and decode a
contradiction. At least \(|\mathcal A|=\binom{k}{d}\) bits are required;
the membership vector suffices for this declared family and its
neighbors. This is not a lower bound on all physically possible
plants, classes with a fixed small number of branches, interactive
reinspection, or the time to decide after reading the body.
If committed source invariants establish persistence for affected
actions as in
\href{https://github.com/eternal-november/eos-research-programmes/pull/618}{the
native eligibility test}, weak justice can suffice and this family
does not justify adopting a more expensive interface.

\paragraph{Checked scope.}
The adapter agrees with direct SCC pruning and an independent
enumeration of recurrent edge supports on 3,500 small authored graphs.
The middle layer check exercises 416 context pairs over \(k=2,3,4\),
with respectively 4, 8 and 64 distinct behavior classes and identical
weak summaries per \(k\). Explicit controls separate prefix filtering,
terminal-only occurrences and unfair stranded interiors. Three Lean
lemmas check equal-rank containment, decoder injectivity and the bit
bound. The graph argument and source-compatibility hypotheses remain
outside the Lean kernel. This is one fixed-policy L2 interface study,
not an independently tested source, risk result, architecture acceptance,
new safety mechanism, or general synthesis proof.
\end{document}
